% OWNER: theory-a
% STATUS: outline
% SOURCES: notes/SPIS_TRESCI_ROBOCZY.md § Rozdz. 2
\chapter{Psychologia środowiskowa i percepcja wnętrz}
\label{ch02:top}
\ChapterStatus{outline}{theory-a}

\section{Teoria środowisk regeneracyjnych}
\label{ch02:sec:art}
\NeedsCite{Attention Restoration Theory; Stress Reduction Theory}
Przegląd ram ART/SRT w kontekście wnętrz mieszkalnych \parencite{kaplan1995restorative}.

\section{Skala postrzeganej regeneracyjności PRS}
\label{ch02:sec:prs}
\Todo{Konstrukt PRS-11 i operacjonalizacja w IDA (ideal/current/target).}
\NeedsCite{hartig1997prs}

\section{Biophilia w projektowaniu wnętrz}
\label{ch02:sec:biophilia}
\Todo{Kellert i implikacje dla parametrów projektowych.}
\NeedsCite{kellert2008biophilic}

\section{Relacja człowiek--środowisko}
\label{ch02:sec:peo}
\Todo{Model PEO / pokrewne ramy; znaczenie dla Room Setup.}

\section{Nastrój, funkcja i kontekst użytkowania}
\label{ch02:sec:nastroj}
\Todo{Mood grid, aktywności, pain points.}

\section{Implikacje dla personalizacji}
\label{ch02:sec:implikacje}
\Todo{Od potrzeb psychologicznych do parametrów projektowych i źródła \texttt{mixed\_functional}.}
